% GENERATED FILE. Edit canonical lessons, then run npm run generate:books.
% canonical-source-hash: fnv1a64:e7a2d2971f802ecc
% canonical-lessons: SA-C538-sakunah, SA-C538-niyatih, SA-C538-mangalam, SA-C538-ganakah, SA-C538-upatisthate

\chapter{Fate and Omens}
\label{ch:sa-fate-and-omens}
\hlchaptermodality{538}

\begin{chapteropening}
\textbf{By the end of this chapter:} \emph{I can talk about omens and what they promise.}

Complete the last lesson of chapter 538: I can talk about omens and what they promise.
\end{chapteropening}

\section[\texorpdfstring{\sk{शकुनः} (śakunaḥ)}{शकुनः (śakunaḥ)}]{\sk{शकुनः} (śakunaḥ) — an omen}

\label{lesson:SA-C538-sakunah}



\begin{quote}
Before the new one: say the Sanskrit for covetous, then the Sanskrit for a gambler.
\end{quote}

\subsection*{You'll want to know: \sk{शकुनः}}
\textbf{\sk{शकुनः}} — \emph{śakunaḥ} — \textquotedblleft{}an omen\textquotedblright{}.

\begin{quote}
\sk{यदि} \sk{शकुनः} \sk{शुभः} \sk{अस्ति}, \sk{तर्हि} \sk{वयं} \sk{श्वः} \sk{यात्रां} \sk{गच्छामः।}
\end{quote}

\begin{quote}
\emph{yadi śakunaḥ śubhaḥ asti, tarhi vayaṃ śvaḥ yātrāṃ gacchāmaḥ.}
\end{quote}

\begin{quote}
If the omen is good, we set out tomorrow.
\end{quote}

A masculine noun.

\begin{cousinweb}
\textbf{\sk{शकुनः}} (\emph{śakunaḥ}) is the word MW gives for a bird, above all a large bird of good or bad omen, and so an omen; MW reports a derivation from √śak, \textquotedblleft{}to be able\textquotedblright{}, only as \textquotedblleft{}said to be\textquotedblright{}, so its origin is uncertain (MW).
\end{cousinweb}

\begin{grammarlens}[title={What you've built}]
One more word for saying what happens if.
\end{grammarlens}

\subsection*{Your turn}
\begin{itemize}
\raggedright
  \item \emph{Say it:} \emph{śakunaḥ}
  \item \emph{Say it:} \emph{śakunaḥ}, once more
  \item \emph{Say it:} \emph{yadi śakunaḥ śubhaḥ asti, tarhi vayaṃ śvaḥ yātrāṃ gacchāmaḥ.} — If the omen is good, we set out tomorrow.
  \item \emph{Recall:} say the Sanskrit for covetous, then the Sanskrit for a gambler, then say \emph{śakunaḥ} again
\end{itemize}

\begin{tcolorbox}[breakable,colback=teal!4,colframe=teal!35!black,title={Before you move on}]
What does \sk{शकुनः} mean? (\textquotedblleft{}An omen\textquotedblright{}.) Say it once more.
\end{tcolorbox}
\section[\texorpdfstring{\sk{नियतिः} (niyatiḥ)}{नियतिः (niyatiḥ)}]{\sk{नियतिः} (niyatiḥ) — destiny}

\label{lesson:SA-C538-niyatih}



\begin{quote}
Before the new one: say the Sanskrit for a gambler, then the Sanskrit for an omen.
\end{quote}

\subsection*{You'll want to know: \sk{नियतिः}}
\textbf{\sk{नियतिः}} — \emph{niyatiḥ} — \textquotedblleft{}destiny\textquotedblright{}.

\begin{quote}
\sk{यदि} \sk{नियतिः} \sk{सत्यम्} \sk{अस्ति}, \sk{तर्हि} \sk{अहं} \sk{किमर्थं} \sk{चिन्तयामि}?
\end{quote}

\begin{quote}
\emph{yadi niyatiḥ satyam asti, tarhi ahaṃ kimarthaṃ cintayāmi?}
\end{quote}

\begin{quote}
If destiny is real, why do I worry?
\end{quote}

A feminine noun.

\begin{cousinweb}
\textbf{\sk{नियतिः}} (\emph{niyatiḥ}) is built on \emph{ni}, \textquotedblleft{}down\textquotedblright{}, and the root √yam, \textquotedblleft{}to hold, to restrain\textquotedblright{}; MW gives \emph{niyati} as the fixed order of things, necessity, destiny (MW).
\end{cousinweb}

\begin{grammarlens}[title={What you've built}]
One more word for saying what happens if.
\end{grammarlens}

\subsection*{Your turn}
\begin{itemize}
\raggedright
  \item \emph{Say it:} \emph{niyatiḥ}
  \item \emph{Say it:} \emph{niyatiḥ}, once more
  \item \emph{Say it:} \emph{yadi niyatiḥ satyam asti, tarhi ahaṃ kimarthaṃ cintayāmi?} — If destiny is real, why do I worry?
  \item \emph{Recall:} say the Sanskrit for a gambler, then the Sanskrit for an omen, then say \emph{niyatiḥ} again
\end{itemize}

\begin{tcolorbox}[breakable,colback=teal!4,colframe=teal!35!black,title={Before you move on}]
What does \sk{नियतिः} mean? (\textquotedblleft{}Destiny\textquotedblright{}.) Say it once more.
\end{tcolorbox}
\section[\texorpdfstring{\sk{मंगलम्} (maṅgalam)}{मंगलम् (maṅgalam)}]{\sk{मंगलम्} (maṅgalam) — auspiciousness; good fortune}

\label{lesson:SA-C538-mangalam}



\begin{quote}
Before the new one: say the Sanskrit for an omen, then the Sanskrit for destiny.
\end{quote}

\subsection*{You'll want to know: \sk{मंगलम्}}
\textbf{\sk{मंगलम्}} — \emph{maṅgalam} — \textquotedblleft{}auspiciousness; good fortune\textquotedblright{}.

\begin{quote}
\sk{यदि} \sk{गृहे} \sk{मंगलम्} \sk{अस्ति}, \sk{तर्हि} \sk{माता} \sk{प्रसन्ना} \sk{भवति।}
\end{quote}

\begin{quote}
\emph{yadi gṛhe maṅgalam asti, tarhi mātā prasannā bhavati.}
\end{quote}

\begin{quote}
If all is well at home, my mother is happy.
\end{quote}

A neuter noun.

\begin{cousinweb}
\textbf{\sk{मंगलम्}} (\emph{maṅgalam}) is the word MW gives for happiness, welfare, and anything auspicious, a good omen; MW reports a derivation from a root √maṅg only after the grammarians, so its origin is uncertain (MW).
\end{cousinweb}

\begin{grammarlens}[title={What you've built}]
One more word for saying what happens if.
\end{grammarlens}

\subsection*{Your turn}
\begin{itemize}
\raggedright
  \item \emph{Say it:} \emph{maṅgalam}
  \item \emph{Say it:} \emph{maṅgalam}, once more
  \item \emph{Say it:} \emph{yadi gṛhe maṅgalam asti, tarhi mātā prasannā bhavati.} — If all is well at home, my mother is happy.
  \item \emph{Recall:} say the Sanskrit for an omen, then the Sanskrit for destiny, then say \emph{maṅgalam} again
\end{itemize}

\begin{tcolorbox}[breakable,colback=teal!4,colframe=teal!35!black,title={Before you move on}]
What does \sk{मंगलम्} mean? (\textquotedblleft{}Auspiciousness; good fortune\textquotedblright{}.) Say it once more.
\end{tcolorbox}
\section[\texorpdfstring{\sk{गणकः} (gaṇakaḥ)}{गणकः (gaṇakaḥ)}]{\sk{गणकः} (gaṇakaḥ) — an astrologer}

\label{lesson:SA-C538-ganakah}



\begin{quote}
Before the new one: say the Sanskrit for destiny, then the Sanskrit for auspiciousness.
\end{quote}

\subsection*{You'll want to know: \sk{गणकः}}
\textbf{\sk{गणकः}} — \emph{gaṇakaḥ} — \textquotedblleft{}an astrologer\textquotedblright{}.

\begin{quote}
\sk{यदि} \sk{गणकः} \sk{सत्यं} \sk{वदति}, \sk{तर्हि} \sk{श्वः} \sk{वर्षा} \sk{भविष्यति।}
\end{quote}

\begin{quote}
\emph{yadi gaṇakaḥ satyaṃ vadati, tarhi śvaḥ varṣā bhaviṣyati.}
\end{quote}

\begin{quote}
If the astrologer is right, it will rain tomorrow.
\end{quote}

A masculine noun.

\begin{cousinweb}
\textbf{\sk{गणकः}} (\emph{gaṇakaḥ}) comes from the root √gaṇ, \textquotedblleft{}to count\textquotedblright{}, which also gives \emph{gaṇayati}, with the agent ending \emph{-aka}; MW gives \emph{gaṇaka} as one who calculates nativities, an astrologer (MW).
\end{cousinweb}

\begin{grammarlens}[title={What you've built}]
One more word for saying what happens if.
\end{grammarlens}

\subsection*{Your turn}
\begin{itemize}
\raggedright
  \item \emph{Say it:} \emph{gaṇakaḥ}
  \item \emph{Say it:} \emph{gaṇakaḥ}, once more
  \item \emph{Say it:} \emph{yadi gaṇakaḥ satyaṃ vadati, tarhi śvaḥ varṣā bhaviṣyati.} — If the astrologer is right, it will rain tomorrow.
  \item \emph{Recall:} say the Sanskrit for destiny, then the Sanskrit for auspiciousness, then say \emph{gaṇakaḥ} again
\end{itemize}

\begin{tcolorbox}[breakable,colback=teal!4,colframe=teal!35!black,title={Before you move on}]
What does \sk{गणकः} mean? (\textquotedblleft{}An astrologer\textquotedblright{}.) Say it once more.
\end{tcolorbox}
\section[\texorpdfstring{\sk{उपतिष्ठते} (upatiṣṭhate)}{उपतिष्ठते (upatiṣṭhate)}]{\sk{उपतिष्ठते} (upatiṣṭhate) — draws near; presents itself}

\label{lesson:SA-C538-upatisthate}



\begin{quote}
Before the new one: say the Sanskrit for auspiciousness, then the Sanskrit for an astrologer.
\end{quote}

\subsection*{You'll want to know: \sk{उपतिष्ठते}}
\textbf{\sk{उपतिष्ठते}} — \emph{upatiṣṭhate} — \textquotedblleft{}draws near; presents itself\textquotedblright{}.

\begin{quote}
\sk{यदि} \sk{शुभः} \sk{समयः} \sk{उपतिष्ठते}, \sk{तर्हि} \sk{वयं} \sk{यात्रां} \sk{गच्छामः।}
\end{quote}

\begin{quote}
\emph{yadi śubhaḥ samayaḥ upatiṣṭhate, tarhi vayaṃ yātrāṃ gacchāmaḥ.}
\end{quote}

\begin{quote}
If a lucky moment comes, we set out on the journey.
\end{quote}

Past (he …): \sk{सः} \sk{गुरुम्} \sk{उपस्थितवान्}, \emph{saḥ gurum upasthitavān}.

\begin{cousinweb}
\textbf{\sk{उपतिष्ठते}} (\emph{upatiṣṭhate}) is built on \emph{upa}, \textquotedblleft{}near\textquotedblright{}, and the root √sthā, \textquotedblleft{}to stand\textquotedblright{}, which also gives \emph{tiṣṭhati}; MW gives \emph{upasthā} as to stand or place oneself near, to be present, to approach; its past participle is \emph{upasthita} (MW).
\end{cousinweb}

\begin{grammarlens}[title={What you've built}]
One more word for saying what happens if.
\end{grammarlens}

\subsection*{Your turn}
\begin{itemize}
\raggedright
  \item \emph{Say it:} \emph{upatiṣṭhate}
  \item \emph{Say it:} \emph{upatiṣṭhate}, once more
  \item \emph{Say it:} \emph{yadi śubhaḥ samayaḥ upatiṣṭhate, tarhi vayaṃ yātrāṃ gacchāmaḥ.} — If a lucky moment comes, we set out on the journey.
  \item \emph{Recall:} say the Sanskrit for auspiciousness, then the Sanskrit for an astrologer, then say \emph{upatiṣṭhate} again
\end{itemize}

\begin{tcolorbox}[breakable,colback=teal!4,colframe=teal!35!black,title={Before you move on}]
What does \sk{उपतिष्ठते} mean? (\textquotedblleft{}Draws near; presents itself\textquotedblright{}.) Say it once more.
\end{tcolorbox}
